% ---------------------------------------------------------------------------
% Squelette de documentation d'une fonction PySpark.
% Fragment insérable dans rapport_stage.tex : pas de préambule, pas de document.
% Supprimez toute section sans contenu réel (pas de « néant », pas de « N/A »).
% ---------------------------------------------------------------------------

\subsection{\texttt{<nom\_fonction>}}
\label{subsec:fn-<nom-fonction-sans-underscore>}

<Une à deux phrases : ce que la fonction produit, à partir de quoi. Pas de mise en
contexte, pas de reformulation de ce que les sections suivantes disent déjà.>

\paragraph{Signature.}\leavevmode\par
\begin{lstlisting}[language=Python]
def <nom_fonction>(<params>) -> DataFrame:
\end{lstlisting}

\paragraph{Arguments.}
\begin{center}
\small
\begin{tabular}{@{}llp{7.6cm}@{}}
\toprule
\textbf{Argument} & \textbf{Type} & \textbf{Rôle} \\
\midrule
\code{<arg>} & \code{<type>} & <rôle ; mentionner la valeur par défaut> \\
\bottomrule
\end{tabular}
\end{center}

% Si la fonction ne prend qu'un ou deux arguments, remplacer le tableau par :
% \begin{itemize}[nosep,leftmargin=*]
%   \item \code{<arg>} (\code{<type>}) --- <rôle>.
% \end{itemize}

\paragraph{Tables lues.}
\begin{itemize}[nosep,leftmargin=*]
  \item \code{<catalog>.<schema>.<table>} --- <ce que la table contient>.
        Colonnes utilisées : \code{<c1>}, \code{<c2>}, \code{<c3>}.
\end{itemize}

% Pour un DataFrame reçu en argument (table amont, pas lue ici) :
% \item \code{<arg>} --- produit par \code{<fonction\_amont>}
%       (\ref{subsec:fn-<...>}). Colonnes attendues : \code{<...>}.

\paragraph{Traitement.}
\begin{itemize}[nosep,leftmargin=*]
  \item <Un bloc logique (variable nommée, colonne calculée) par tiret, une
        phrase courte et précise chacun : ce que le bloc représente + son info
        clé (filtre décisif, type de jointure, clé). Pas un tiret par appel
        Spark ; une expression chaînée simple tient en un ou deux tirets.>
\end{itemize}

\paragraph{Table produite.}
Une ligne par <granularité : produit, session, paire de produits\ldots>.
\begin{center}
\small
\begin{tabular}{@{}llp{6.8cm}@{}}
\toprule
\textbf{Colonne} & \textbf{Type} & \textbf{Provenance et sens} \\
\midrule
\code{<col>} & \code{<type>} & <table.colonne d'origine, ou expression> \\
\bottomrule
\end{tabular}
\end{center}
